\section{Experimental Setup}
\label{sec:experimental-setup}

\subsection{Evaluation Scope}
Synthetic scores assess categorical search; a separate real-model audit assesses recommendation utility and slates. Their score scales and evidence are kept distinct.

The initial synthetic comparison uses 24 attributes with cardinalities $(2,4,3)$ repeated eight times, landscape seed 2026, three searches per method, base search seed 42, population 20 per end, 4,000 requests, and threshold 0.65. All assignments are feasible.

Separately recorded post hoc controls add no model training or test-based tuning. They cross three score families, three seeded instances per family (2026--2028), three encodings (identity and two random relabelings), and five search seeds (1001--1005). The initial extension records five methods including uniform random (675 searches); a separately recorded revision adds ABC and SA under the same configuration (270 searches), totaling 945 searches at 4,000 requests each. Relabeling preserves semantic scores. Encodings and searches are not independent instances; threshold counts are interpreted within family.

Additive scores average independent seeded category weights in $[0,1]$. NK/trap scores combine mean three-attribute lookup values (weight 0.65) with planted triple rewards (weight 0.35): a full pattern match receives one, otherwise a triple with $h$ matches receives $(2-h)/3$. Modular scores count the fraction of eight disjoint triples whose semantic-category sum matches a seeded residue modulo three. Unequal category cardinalities can leave marginal associations, so this family is not claimed to contain only marginally invisible interactions.

\subsection{Dataset and Target Model}
The real-model experiment uses MovieLens-1M ratings, categorical user metadata, and movie genres. Ratings of at least four are treated as positive interactions, and users require at least three positives. The resulting checkpoint contains 6,035 users and 3,706 items.

For each user, the last two chronological positives form validation and test interactions; earlier positives form training data. NeuMF uses embedding dimension 16 and multilayer perceptron widths 64 and 32. Training uses Adam with learning rate $10^{-3}$, batch size 1,024, four sampled negatives per positive, and a 30-epoch cap. Development-user validation selects checkpoints with patience five and minimum improvement $10^{-4}$. For seeds 42, 100, and 2026, training stops after 25, 18, and 13 epochs, retaining checkpoints from epochs 20, 13, and 8, respectively. These are validation-selected checkpoints; the stopping rule does not establish mathematical convergence.

Users are partitioned once into 1,810 development users and 4,225 audit users using split seed 314159. Feature selection and checkpoint selection use development users; subgroup search uses validation interactions of audit users, followed by confirmation on their test interactions. The development/audit partition controls selection and evaluation; the recommender trains on earlier interactions from all eligible users.

Validation ranking excludes training positives; test ranking additionally excludes the validation positive. Evaluation uses full-catalog NDCG@10 with one held-out relevant item per user. Demographic attributes are used for auditing and do not enter the recommender.

\subsection{Baselines and Computational Controls}
The initial comparison evaluates native FairRec DPSO~\cite{guo2023fairrec}, categorical GA/ACO, and random-key DE. Every request counts, including repeats. Challengers share initial candidates and split budgets between maximizing and minimizing ends; native DPSO retains its initializer and coupled policy. Results therefore compare configurations, including initialization, under equal total requests.

The initial NeuMF comparison uses population 20, 4,000 requests per method/seed, support 20, and threshold 0.02. Three paired model/search seeds yield 12 successful runs. Cached search times exclude training and objective setup.

The initial demographic grammar contains gender, age, and occupation. Backward elimination starts from all three attributes and uses three-fold cross-validation and a depth-five decision tree, with at most two retained features. Labels indicate training activity below the development median (58 positives), a proxy rather than a measured fairness outcome. Eliminated attributes remain unconstrained.

For reproduction of the saved runs, the archived experimental code includes mRMR and Chi-square ranking hooks, each configured to retain the top ten attributes before taking their union. Each evaluated real-data schema has three attributes, so these hooks retain the entire schema and perform no reduction. They are disclosed solely as settings of our archived implementation and are omitted from the evaluated method in Section~\ref{sec:feature-selection}. The hooks belong to our experimental implementation; this study does not claim to reproduce a FairRec feature-selection pipeline. No mRMR, Chi-square, or fusion benefit is evaluated or claimed.

The checkpoint extension crosses ten independent search seeds (1001--1010) with each of the three frozen models. Native DPSO, GA, ACO, DE, and uniform random sampling are compared in both the original selected grammar and the unfiltered gender--age--occupation grammar, producing 300 searches. Each run retains the 4,000-request budget, population 20, support 20, and disparity threshold 0.02. Random sampling uses the challengers' initial candidates followed by independent uniform category assignments; repeats and invalid requests count. Exact enumeration separately scores every cell once, establishing within-grammar optima, regret, threshold recall, and optimum-attainment rates. It stops at space exhaustion rather than consuming an artificial 4,000 requests. The unfiltered space contains only 294 cells, so this is a correctness and selection audit rather than a scalability test.

The revision adds ABC and SA for all 60 model--grammar--search-seed combinations, bringing this comparison to 420 searches across seven methods. Earlier result folders are preserved; added indexed cell scores are checked against every original exact score before search. Cross-method runtimes from the two implementations are not pooled.

\subsection{Optimization Hyperparameters}
All heuristic challengers use 20 initial candidates per end and split 4,000 requests equally between maximizing and minimizing signed scores. Shared starts are sampled with NumPy SeedSequence $[s,0]$ and challenger search randomness with $[s,1]$. Cached duplicates, unsupported cells, and initialization consume requests; there is no early stopping or parameter tuning against test outcomes.

GA uses size-three tournaments, uniform crossover with probability 0.5 per coordinate, mutation probability $1/d$ per coordinate to a different category, and one cached elite. ACO starts trails at one, evaporates 20\% per batch, reinforces the best $\lfloor20/3\rfloor=6$ valid ants by $1/(r+1)$ for zero-based rank $r$, clips trails to $[0.05,20]$, and mixes trail sampling with 10\% uniform exploration; its heuristic is constant. DE uses DE/rand/1/bin random keys, differential weight 0.7, crossover 0.8 with one forced donor key, clipping to $[0,1]$, argmax decoding, and nonworsening replacement. ABC uses 20 employed and 20 rank-weighted onlooker attempts per complete cycle, one-coordinate donor moves, strict-improvement replacement, and abandonment after ten failed attempts. SA starts at the best initialized candidate per end, changes one coordinate, and cools geometrically from 0.15 to 0.001 over the remaining requests with Metropolis acceptance.

Native DPSO retains thermal coefficient 0.09, cognitive/social coefficients 2.0, velocities clipped to $[-2,2]$, integer position conversion, and category-bound clipping. Its uniform category initializer and coupled two-ended control flow remain native; equal budgets do not imply identical trajectories or starts. The generation cap is 4,000, with the request ledger enforcing the earlier total-request stop. Exhaustive evaluation ends at grammar exhaustion.

\subsection{Feature-Selection Ablation Design}
The MovieLens target-by-cap ablation crosses activity-proxy versus development-utility-deficit labels with maximum subset sizes two and three, retaining the unfiltered reference. At each checkpoint, exact scoring evaluates all five arms at supports 20, 50, 100, and 200 (60 cases). Backward elimination starts from the complete schema and uses three-fold stratified cross-validation, balanced accuracy, a one-hot encoded depth-five decision tree, minimum one feature, and tolerance 0.001. Caps are upper bounds, not required final sizes. This design examines target and wrapper-constraint effects jointly with support sensitivity.

\subsection{Additional Frozen-Model Datasets}
We audit the locally supplied FairRec-processed LFM360K and Amazon tables and frozen DeepFM checkpoints~\cite{guo2023fairrec,guo2017deepfm}. No additional model training is performed. All current model variables are verified as restored; unused optimizer slots are permitted. Embedding dimension is four; inference uses batches of 8,192 with four intra-operation threads and one inter-operation thread under Python 3.8.0, TensorFlow 2.3.0, and DeepCTR 0.9.3. Original checkpoint training seeds, stopping histories, and all preprocessing decisions cannot be reconstructed, so these are one-checkpoint supplementary audits per dataset.

Each supplied table contains one relevant item and 49 sampled negative rows per eligible user. We calculate per-user NDCG@10 on these 50 rows, ordering by descending predicted score and ascending item ID for ties; duplicate candidate rows are retained. These utilities are not comparable in level to MovieLens full-catalog NDCG. The eligible extracts contain 267,034 LFM360K and 175,878 Amazon users. LFM360K uses gender, age, and country. Amazon uses history level, mean-price level, and maximum-price level. These are behavioral profiles, not demographic protected groups; the supplied extract does not establish its original Amazon category/release or leakage-free history construction.

Using split seed 314159 separately within each table, 30\% of users form development, 35\% discovery, and the remainder confirmation, with integer counts rounded down for the first two partitions. The metadata vocabulary is declared from the supplied user table without using utility outcomes. Feature selection uses only development utility-deficit labels, with caps two and three. Discovery compares selected and unfiltered grammars at support 20 and threshold 0.02. Seven methods and ten search seeds give 140 searches per dataset (280 total), at 4,000 requests each. Both datasets have supported cells in both grammars, and exact enumeration supplies within-grammar references. Both caps and no selection are scored at supports 5, 20, 50, and 100 (24 ablation cases across the two datasets). One discovery-maximizing rule per grammar is fixed before confirmation on disjoint users; pointwise bounded-outcome intervals use the formula below. Disjoint confirmation users reduce outcome reuse but do not verify the supplied checkpoint's historical training separation or eliminate shared-model dependence.

\subsection{Projection and Fixed-Rule Confirmation}
To measure attenuation by feature removal, each supported above-threshold full cell is mapped to its selected-feature parent. We count how many full cells have parents that still exceed 0.02, keeping the audit users and global reference fixed; multiple cells may share one parent. This measures aggregation across eliminated attributes, not optimizer regret between equivalent grammars. In MovieLens, one exhaustive validation maximizer per model and grammar is fixed using enumeration order to break ties, then evaluated on test interactions. Exploratory pointwise intervals use 2,000 user-bootstrap resamples, preserving subgroup membership in the global reference. To guard against zero bootstrap variance in sparse zero-hit groups, we also report conservative bounded-outcome intervals. Writing $N$ for audit users, $n$ for subgroup support, and $p=n/N$, the gap is $(1-p)(\overline q_{\neg g}-\overline q_g)$. Two Hoeffding bounds and a union bound give a fixed-rule 95\% radius
\begin{equation}
(1-p)\left[\sqrt{\frac{\log 80}{2n}}+\sqrt{\frac{\log 80}{2(N-n)}}\right].
\end{equation}
This calculation assumes independent bounded user outcomes and fixes the observed group proportion. In MovieLens, neither interval adjusts for adaptive selection or multiplicity; shared audit users and post hoc design prevent interpreting them as independent population confirmation. Additional-dataset bounds use disjoint confirmation users but remain pointwise and conditional on the supplied checkpoint.

\subsection{Measurements}
We report attained maxima, first attainment, unique assignments, above-threshold breadth, invalid requests, and runtime. Progress uses objective requests rather than incompatible native generation counts.

Calibration and exposure-opportunity measurements accompany utility-selected groups without modifying fitness. These descriptive measurements complement the distinct explanatory and intervention-based questions addressed by counterfactual approaches~\cite{li2021towards,medda2024gnnuers}. The opportunity distribution is explicitly declared because exposure comparisons depend on the selected reference~\cite{singh2018fairness}.

Extension timings include cached objective scoring, optimization, and in-memory request recording; frozen-model inference and compressed trace output are excluded. They therefore differ from the original expanded-audit timings. Extension traces record every request, including duplicates and unsupported candidates. Category-permutation summaries report changes across recorded encodings; because encoded initial candidates are shared, rather than semantically coupled across encodings, these changes do not isolate numerical movement from initialization and finite-search randomness.
